\documentclass[11pt]{article}
\usepackage[T1]{fontenc}
\usepackage{lmodern}
\usepackage{amsmath,amssymb,amsthm,mathtools}
\usepackage[a4paper,margin=29mm]{geometry}
\usepackage{microtype}
\usepackage{booktabs,longtable,array}
\usepackage{enumitem}
\usepackage[hidelinks]{hyperref}
\usepackage{url}
\setlength{\emergencystretch}{2em}
\setlist{itemsep=3pt,topsep=5pt}
\newtheorem{theorem}{Theorem}[section]
\newtheorem{proposition}[theorem]{Proposition}
\newtheorem{lemma}[theorem]{Lemma}
\newtheorem{corollary}[theorem]{Corollary}
\theoremstyle{definition}
\newtheorem{definition}[theorem]{Definition}
\newtheorem{conjecture}[theorem]{Conjecture}
\theoremstyle{remark}
\newtheorem{remark}[theorem]{Remark}
\newcommand{\PP}{\mathcal P}
\newcommand{\Cl}{\mathcal C}
\newcommand{\cT}{\mathcal T}
\newcommand{\sig}{\sigma}
\newcommand{\ee}{\mathrm e}
\newcommand{\LPF}{P^{+}}
\newcommand{\card}[1]{\lvert #1\rvert}
\newcommand{\norm}[1]{\lVert #1\rVert}
\newcommand{\floor}[1]{\left\lfloor #1\right\rfloor}
\newcommand{\lean}[1]{\nolinkurl{#1}}
\DeclareMathOperator{\lcm}{lcm}
\title{Density amplification in a largest-prime-factor dynamical system}
\author{Jonas Whidden}
\date{15 October 2026\\\small Research draft}
\begin{document}
\maketitle

\begin{abstract}
For a finite set of primes \(S\), retain its elements and adjoin
\(P^{+}(2(p+q)-1)\) for every \(p,q\in S\), allowing \(p=q\) at every
generation. We prove that there are constants \(\eta>0\) and \(M_0\)
such that any finite seed missing at most \(\eta M/\log M\) primes
up to one scale \(M\ge M_0\) generates a set containing all but
\(O_{S,M}(X^{79/80})\) primes up to \(X\). A stronger statement controls
the actual finite generations at geometrically increasing cutoffs.
The proof couples a power-saving exceptional-set estimate for canonical
prime-pair representations with a deletion estimate uniform over all
missing input sets. Exact representation-degree normalization and an
averaged singular-factor overlap bound give a defect multiplier \(17/16\),
smaller than the available scale growth. We give the complete iteration
and interpolation argument, identify the analytic inputs and their
normalizations, and record the corresponding Lean declarations.
The motivating conjectures assert that the seed \(\{2\}\) generates every
prime, and that each odd-prime singleton generates every odd prime.
The theorem applies to sufficiently rich finite seeds; neither conjecture
is proved here.
\end{abstract}

\section{Introduction}

Largest-prime-factor dynamics combines two operations with rather different
behavior. Addition produces many values from a finite set; taking the largest
prime factor merges many of those values into a single output. A large number
of pairs therefore need not give a large number of distinct new primes.
This paper studies an accumulating, branching system in which that collision
problem can be controlled after a suitable finite density condition has
been reached.

Write \(P^{+}(n)\) for the largest prime factor of \(n>1\). For a finite
set \(S\) of primes define
\begin{equation}\label{eq:step}
 \cT(S)=S\cup\{P^{+}(2(p+q)-1):p,q\in S\},
 \qquad S_0=S,\quad S_{j+1}=\cT(S_j).
\end{equation}
Pairs are ordered when counted, although the set of outputs is unchanged
by interchanging the two inputs. Diagonal pairs are always permitted.
The generated closure is \(\Cl(S)=\bigcup_{j\ge0}S_j\).

Our main result is a robust propagation theorem: a sufficiently dense
finite prime prefix amplifies into density one, with a quantitative
power-saving exceptional set. The robustness is important. We impose
no independence assumption, no random model, and no residue-distribution
hypothesis on the missing primes of the seed. Their effects are controlled
jointly through an incidence matrix.

\subsection{The motivating generation conjectures}\label{sec:conjectures}

Let \(\PP\) denote the set of all primes. The original objective is
complete prime generation, not merely relative density one.

\begin{conjecture}[Generation from two]\label{conj:two}
For the dynamics \eqref{eq:step},
\[
 \Cl(\{2\})=\PP.
\]
Equivalently, for every prime \(r\) there is a finite generation \(j\)
with \(r\in S_j\) when \(S_0=\{2\}\).
\end{conjecture}

\begin{conjecture}[Generation from an odd prime]\label{conj:odd}
For every odd prime \(p\),
\[
 \Cl(\{p\})=\PP\setminus\{2\}.
\]
The generation at which a target prime appears may depend on both the
initial prime \(p\) and the target. Repeated inputs remain permitted at
every generation.
\end{conjecture}

The missing prime \(2\) in Conjecture~\ref{conj:odd} is forced: every
new output is odd. The odd-singleton conjecture also implies
Conjecture~\ref{conj:two}. Indeed \(P^{+}(2(2+2)-1)=7\), and monotonicity
of the closure gives
\(\Cl(\{7\})\subseteq\Cl(\{2\})\).
If Conjecture~\ref{conj:odd} holds even just for \(p=7\), this inclusion
supplies all odd primes, while \(2\) is present initially.

The density-amplification theorem below is a step toward these
generation questions, rather than their replacement. Even proving dense
entry from a singleton would leave a further gap: the resulting
\(O(X^{79/80})\) exceptional set may contain infinitely many primes.
Complete generation requires showing that each individual prime is
eventually reached, including any primes left outside that density
conclusion.

\subsection{Related work and the scope of the contribution}

Caragiu and Vicol studied prime magmas with operation
\(P^{+}(ap+bq)\), including accumulation from all available pairs
\cite{CV}. Their cyclicity questions and finite-generation examples are
close in spirit, but their homogeneous operations are different from
the affine operation in \eqref{eq:step}. Back and Caragiu considered
greatest-prime-factor recurrences with two preceding inputs
\cite{BC}; these are single trajectories, rather than branching
closures. Neither comparison identifies an existing theorem about the
specific closure studied here.

The analytic setting also has substantial antecedents. Ramar\'e and
Ruzsa developed additive results for dense subsets of sifted sequences
using sieve methods \cite{RR}. Restriction estimates for sieve
majorants provide another approach to additive questions in the primes
\cite{GT}. Alsetri and Shao showed that even prime subsets with arbitrarily
high global relative density can miss a positive proportion of even
sums; they also established an almost-all binary Goldbach theorem under
uniform relative-density conditions in reduced residue classes
\cite{AS}. These results explain why a density-to-representation passage
requires care. In particular, an unspecified \(o(X)\) exceptional set
among integers is insufficient when the desired targets are indexed by
only \(O(X/\log X)\) primes.

The contribution here is the combination of canonical largest-prime
ownership, a deletion estimate uniform over the actual missing input
set, and iteration on the \emph{generated states}. The exponent \(79/80\)
is inherited from a deliberately conservative analytic supply estimate;
it is not asserted to be optimal. The character estimates used by the
installed proof chain draw on the 2026 zero-free work of OpenAI
\cite{OAI}. We distinguish that analytic source from the dynamical
argument and state the exact estimate consumed below. We make no claim
to a new zero-free region, a new Goldbach exponent, or a resolution of
the small-seed reachability problem.

\section{An informal explanation of the argument}\label{sec:informal}

Imagine that the current set already contains almost every prime below a
large number \(M\). We try to produce almost every prime in the next window,
from \(M\) to roughly \(5M/4\). A target prime \(r\) will be produced if we
can find two available primes whose sum is a carefully chosen even number.

There are two choices for that even number. If \(r\) is \(3\) modulo \(4\),
use \((r+1)/2\). Then the value before taking its largest prime factor is
exactly \(r\). If \(r\) is \(1\) modulo \(4\), use \((3r+1)/2\). The value
is then \(3r\), whose largest prime factor is again \(r\). Thus the
multiplicative output question becomes a controlled pair of additive
representation questions.

Almost every target has many representations using the full prime prefix.
Some of those representations can disappear because an input prime is
missing from the current state. A crude argument charging each missing
input separately loses too much. Instead we examine all failed targets
together. If a target fails, every representing pair must touch a missing
input. Cauchy--Schwarz converts that statement into an energy inequality.
The diagonal part of this energy is known exactly after dividing each
target row by its own representation count. The overlap between two
different rows is bounded by a three-prime sieve and then averaged over
the actual failed targets.

This gives the key comparison. The number of holes at the next cutoff is
at most \(17/16\) times the old number, plus a power-saving exceptional
term and a logarithmic remainder. Meanwhile the cutoff grows by roughly
\(5/4\). The holes therefore occupy a smaller and smaller fraction of the
prime prefix. Once this vanishing fraction is available, the stronger
power-saving recurrence can be used permanently. Its forcing term has
size \(M^{79/80}\); geometric summation gives the same exponent for the
holes themselves.

Finally, a prime that is missing forever is missing at every finite
generation. Choosing the first cutoff above an arbitrary \(X\) transfers
the finite-generation bound to the permanently missing primes below
\(X\). This last step matters: a bound only at a sparse sequence of
cutoffs is not yet the advertised statement for every large \(X\).

The argument begins \emph{after dense entry}. The fact that a small seed
has many pairs, or behaves well modulo a few small primes, does not prove
that it reaches this regime. The seed \(\{2\}\) remains the main unresolved
instance for this method. Our ultimate target is the complete generation
assertion in Conjecture~\ref{conj:two}, together with its odd-singleton
counterpart, Conjecture~\ref{conj:odd}.

\section{Definitions and main results}

Let
\[
 \PP(M)=\{p\le M:p\text{ is prime}\},\qquad
 W(M)=\PP(\floor{5M/4})\setminus\PP(M),
 \qquad N(M)=\frac{M}{\log M}.
\]
All asymptotic statements concern positive integer \(M\) tending to
infinity. Floors in cutoffs are retained. Define
\begin{equation}\label{eq:cutoffs}
 m_0=M,\qquad m_{j+1}=\floor{5m_j/4},\qquad
 E_j=\card{\PP(m_j)\setminus S_j},\qquad \gamma=\frac{79}{80}.
\end{equation}
The iterated cutoffs \(m_j\) are distinct from the fixed entry threshold.

\begin{theorem}[Power-saving propagation]\label{thm:main}
There exist \(\eta>0\) and an integer \(M_{\mathrm{entry}}\ge8\) such that,
for every \(M\ge M_{\mathrm{entry}}\) and every finite set \(S\) of primes
satisfying
\begin{equation}\label{eq:entry}
 \card{\PP(M)\setminus S}\le\eta\frac{M}{\log M},
\end{equation}
there are \(K>0\) and an integer \(J\) for which
\begin{equation}\label{eq:stage}
 \card{\PP(m_j)\setminus S_j}\le K m_j^{79/80}
 \qquad(j\ge J).
\end{equation}
In particular, for some \(K'>0\) and \(X_0\),
\begin{equation}\label{eq:closure}
 \card{\{r\le X:r\text{ prime},\ r\notin\Cl(S)\}}
 \le K'X^{79/80}\qquad(X\ge X_0).
\end{equation}
The constants \(K,J,K',X_0\) may depend on the fixed seed and entry scale.
The entry fraction and entry threshold precede the quantifiers over
\(M\) and \(S\).
\end{theorem}

In fact the Lean statements allow an arbitrary finite set of natural
numbers as the initial set: the defect counts only missing primes.
We state the prime-seed specialization to keep the dynamical
interpretation uniform.

\begin{corollary}\label{cor:density}
For every seed satisfying Theorem~\ref{thm:main},
\[
 \frac{\card{\Cl(S)\cap\PP(X)}}{\pi(X)}\longrightarrow1.
\]
Consequently \(\Cl(S)\) is infinite. The fraction of primes below \(X\)
that remain unreachable is \(O_{S,M}(X^{-1/80}\log X)\).
\end{corollary}

\begin{proof}
Divide \eqref{eq:closure} by \(\pi(X)\sim X/\log X\).
Since \(X^{-1/80}\log X\to0\), the missing fraction tends to zero.
An infinite number of primes is therefore reached.
\end{proof}

\begin{corollary}[Finite-stage appearance]\label{cor:time}
Under the same hypotheses, all but \(O_{S,M}(X^{79/80})\) primes up to
\(X\) are present after \(O_M(\log X)\) generations. The lower medians
of the increasing finite prime sets \(S_j\) tend to infinity.
\end{corollary}

\begin{proof}
For \(m\ge8\), \(\floor{5m/4}\ge9m/8\). The first index \(j\)
with \(m_j\ge X\) is therefore at most
\(\lceil\log(X/M)/\log(9/8)\rceil\) when \(X>M\).
For sufficiently large \(X\), that index is at least \(J\), and
\(m_j\le5X/4\). Apply \eqref{eq:stage}. To prove the median assertion,
fix \(B\). There are at most \(\pi(B)\) primes at most \(B\).
The sets \(S_j\) are nested and their union is infinite, so their
cardinalities tend to infinity. Eventually more than half their elements
exceed \(B\).
\end{proof}

The two corollaries above are ordinary mathematical consequences of the
formalized main theorem and the prime number theorem. They are not
advertised as separate Lean declarations in this release.

\section{Canonical prime targets and exact deletion}

For an odd prime \(r\), put
\begin{equation}\label{eq:targets}
 c_r=\begin{cases}1&r\equiv3\pmod4,\\3&r\equiv1\pmod4,\end{cases}
 \qquad t_r=\frac{c_rr+1}{2}.
\end{equation}
Then \(t_r\) is even and \(2t_r-1=c_rr\). Since every prime factor of
\(c_r\) is at most \(r\), we have \(P^{+}(c_rr)=r\). Therefore
\begin{equation}\label{eq:target-production}
 p+q=t_r,\quad p,q\in S\quad\Longrightarrow\quad r\in\cT(S).
\end{equation}
This includes \(p=q\).

For a finite prime set \(P\) and an even integer \(t\), define the
\emph{ordered} representation degree
\[
 R_P(t)=\card{\{p\in P:t-p\in P\}},
\]
where subtraction is interpreted only when \(p\le t\).
Every representation is counted by its first input; a diagonal
representation contributes once. Write \(R_P(r)=R_P(t_r)\) when the
target prime is clear. For \(D=P\setminus S\), the absence of a
representation with both inputs in \(S\) implies
\begin{equation}\label{eq:half-deleted}
 R_P(r)\le2\card{\{p\in D:t_r-p\in P\}}.
\end{equation}
Indeed the involution \(p\mapsto t_r-p\) pairs the representations;
every pair meets \(D\). The same inequality is valid at a fixed point
of the involution, so no diagonal correction is discarded.

\begin{lemma}[Distinct targets]\label{lem:injective}
For \(M\ge3\), the map \(r\mapsto t_r\) is injective on \(W(M)\).
\end{lemma}
\begin{proof}
Within either residue branch this is immediate. An equality across
branches would give \(r=3s\) or \(s=3r\), whereas both targets lie in
\((M,5M/4]\). Their ratio is less than \(5/4\), so this is impossible.
\end{proof}

We use the finite divisor factor
\begin{equation}\label{eq:sigma}
 \sig(n)=\prod_{\substack{\ell\mid n\\\ell\text{ prime},\,\ell\ge3}}
       \left(1+\frac1{\ell-2}\right),\qquad \sig(1)=1.
\end{equation}
Only nonzero differences occur when \(\sig(|t_r-t_s|)\) is used.
The two normalizations in the proof have different roles:
\(\sig(t_r)^{-1}\) removes the arithmetic factor from a supply or sieve
bound, while \(R_P(r)^{-1}\) makes each target row have mass exactly one.

\section{The analytic representation supply}\label{sec:source}

We state the quantitative inputs with their ranges and normalizations.
They are proved in the maintained Lean source, rather than introduced
as new assumptions of Theorem~\ref{thm:main}.

\subsection{Character input and the complete error}

The centered character estimate consumed by the major-arc calculation is
\begin{equation}\label{eq:char}
 \left|\sum_{n\le x}\Lambda(n)\chi(n)
       -x\,\mathbf1_{\chi=\chi_0}\right|
 \le A x^{15/16}
       (1+\log q+\log(x+2))^2
 \qquad(x\ge65536).
\end{equation}
Here one constant \(A>0\) precedes every modulus \(q\ge1\), character
\(\chi\bmod q\), and prefix \(x\). The principal character includes the
finite Euler-factor correction in the proof. The nonprincipal branch
uses the installed \(7/8\) character estimate, while the principal
branch is relaxed to \(15/16\). The formal owner is
\lean{QRH.exists_all_characterChebyshevSum_power_bound}.
The zero-free theorem in \cite{OAI} is part of the analytic provenance;
we use the quantitative estimate \eqref{eq:char}, not a bare
zero-free statement with an unstated uniformity inference.

For an integer interval \(I=(s,s+L]\), let
\[
 V_I(t)=\sum_{\substack{a,b\in I\\a+b=t}}\Lambda(a)\Lambda(b),
 \qquad
 H_I(t)=\card{\{(a,b)\in I^2:a+b=t\}}.
\]
Use the standard binary Goldbach singular series
\[
 \mathfrak S(t)=
 \begin{cases}
 2\displaystyle\prod_{\ell>2}\left(1-\frac1{(\ell-1)^2}\right)\sig(t),
       &t\text{ even},\\
 0,&t\text{ odd}.
 \end{cases}
\]
The error is centered at the \emph{infinite} series:
\(\mathcal E_I(t)=V_I(t)-H_I(t)\mathfrak S(t)\).

\begin{lemma}[Complete fixed-interval error]\label{lem:source-error}
There is \(A_1>0\) such that, for every integer \(M\ge2^{18}\), every
interval with
\[
 M/4\le s,\quad s+L\le2M,\quad L\le M,
\]
and every \(T\subseteq\{1,\ldots,2M\}\),
\begin{equation}\label{eq:source-error}
 \sum_{t\in T}|\mathcal E_I(t)|^2\le A_1 M^{239/80}.
\end{equation}
The constant is uniform in the interval and target set.
\end{lemma}

Here is the error accounting behind this bound. Choose an integer
\(P\) with
\(\tfrac12M^{1/40}\le P\le M^{1/40}\).
The cyclic grid in the formal proof has a period \(N>2M\), divisible
by \(M\) and every \(q\le P\); one may enlarge such a period to ensure
\(s+L<N\). Consequently neither the interval nor the target set wraps
around the grid. Inverse Fourier inversion gives the exact identity
\begin{equation}\label{eq:four}
 \mathcal E_I(t)=\mathcal E_{\mathrm{maj}}(t)
   +\mathcal E_{\mathrm{min}}(t)+\mathcal E_{\mathrm{ker}}(t)
   +H_I(t)(\mathfrak S_P(t)-\mathfrak S(t)).
\end{equation}
The kernel packet retains the interval weights; the last packet retains
the actual pair count. Neither is replaced by an unweighted tail.

For completeness, the squared packet budgets used in the proof, with
\(L_M=1+\log(2M)\), are bounded by constants times
\begin{align}
 &(8K^2(1+\log M)^4+6K^4(1+\log M)^8)M^{119/40},
                                      &&\text{major},\label{eq:maj}\\
 &C M^{1/160}(M^3/P+M^{13/5})L_M^9,
                                      &&\text{minor},\label{eq:min}\\
 &B M^{481/160}/P^2,                   &&\text{kernel},\label{eq:ker}\\
 &D M^{481/160}/P,                     &&\text{series tail}.\label{eq:tail}
\end{align}
All four constants are positive and independent of \(M,N,I,T\).
The major bound uses \eqref{eq:char} after the character expansion at
reduced rational frequencies. The minor bound uses the installed
Vaughan decomposition and its mean-square estimate. The last two
bounds use Fourier inversion for the interval kernel and divisor
moment bounds for the singular tail. The exact finite Fourier
identities and the four estimates have separate formal owners listed
in Appendix~\ref{app:formal}; thus an overlap or kernel error is not
hidden inside an unspecified error term.

\begin{proof}[Proof of Lemma~\ref{lem:source-error} from the packet bounds]
Apply \(|u_1+\cdots+u_4|^2\le4\sum_i|u_i|^2\) to
\eqref{eq:four} and sum over the same target set.
Substituting \(P\asymp M^{1/40}\) in \eqref{eq:maj}--\eqref{eq:tail}
bounds the result by
\[
 A_2 M^{477/160}(1+\log(2M))^9.
\]
For example, the minor term \(M^{1/160}M^3/P\) has exponent
\(477/160\), the kernel term has exponent \(473/160\), and the series
tail again has exponent \(477/160\). The other terms are smaller.
The elementary uniform bound
\[
 (1+\log(2M))^9\le1442^9 M^{1/160}\qquad(M\ge1)
\]
absorbs the remaining logarithms. Since \(478/160=239/80\),
\eqref{eq:source-error} follows. The formal proof constructs \(P\)
by an integer cutoff and includes the factor-two rounding bounds.
\end{proof}

The packet estimates are analytic library results, with their own proofs
in the companion formal development. The elementary argument above
records their complete consumer; it does not purport to reprove the
underlying character and minor-arc theorems from first principles.

\subsection{Supply on the two canonical branches}

\begin{proposition}[Positive supply]\label{prop:supply}
There is \(C_{\mathrm{src}}>0\) such that, for every \(M\ge2^{44}\),
there is \(Z_M\subseteq W(M)\) with
\[
 \card{Z_M}\le C_{\mathrm{src}}M^{79/80}
\]
and, for every \(r\in W(M)\setminus Z_M\),
\begin{equation}\label{eq:supply}
 \frac{R_{\PP(M)}(r)}{\sig(t_r)}
 \ge a\frac{M}{\log^2 M},\qquad a=\frac{63}{1024}.
\end{equation}
\end{proposition}

\begin{proof}
For the branch \(r\equiv1\pmod4\), use the fixed input interval
\((\lceil M/4\rceil,M]\). For \(r\equiv3\pmod4\), use the analogous
interval at scale \(\floor{M/2}\). Both inputs are at most \(M\).
The interval-pair geometry, with all integer endpoints retained, gives
the positive comparison underlying
\[
 V_I(t_r)\ge \frac1{16}M\sig(t_r)
\]
outside an exceptional set. To see the size of that set, a failure of
this lower bound gives a fixed positive multiple of \(M\sig(t_r)\)
in the error \(\mathcal E_I(t_r)\). Since \(\sig(t_r)\ge1\),
\eqref{eq:source-error} and Chebyshev's inequality bound the number
of failures by
\[
 O(M^{239/80}/M^2)=O(M^{79/80}).
\]
The map in Lemma~\ref{lem:injective} prevents a multiplicity loss when
passing from even targets to primes. The two interval branches are
counted together in the constant.

The weights \(\Lambda(a)\Lambda(b)\) include prime powers. Their
total unwanted contribution is at most
\(4\sqrt M\log^2 M\); the explicit comparison at \(M\ge2^{44}\)
bounds this by \(M/1024\). After discarding prime powers, each
remaining weight is at most \(\log^2M\).
Using \(\sig(t_r)\ge1\), the remaining weighted supply is at least
\[
 \left(\frac1{16}-\frac1{1024}\right)M\sig(t_r)
   =\frac{63}{1024}M\sig(t_r).
\]
Division by \(\log^2M\) proves \eqref{eq:supply}.
The finite interval geometry and exceptional-set comparison are
formalized in \lean{LPFWideSource}; the prime-power subtraction is
performed before the conversion to ordered prime representations.
\end{proof}

The endpoint \(2^{44}\) is a source-conversion cutoff. It is not the
entry threshold of Theorem~\ref{thm:main}.

\section{Uniform deletion control}\label{sec:deletion}

Put \(P=\PP(M)\), \(D=P\setminus S\), \(E=\card D\), and
\(G=W(M)\setminus Z_M\). For \(r\in G\), let
\[
 w_r(p)=\frac{\mathbf1_{\{p\in P,\ t_r-p\in P\}}}{R_P(r)}.
\]
By \eqref{eq:supply}, every denominator is positive. The row sums
are exactly one, and their diagonal energies satisfy
\begin{equation}\label{eq:self}
 \sum_{p\in P}w_r(p)=1,\qquad
 \sum_{p\in P}w_r(p)^2=\frac1{R_P(r)}
       \le U,\qquad U=\frac{\log^2M}{aM}.
\end{equation}

The three-prime upper sieve and exact normalization give, for distinct
\(r,s\in G\),
\begin{equation}\label{eq:cross}
 \sum_{p\in P}w_r(p)w_s(p)
 \le V\sig(|t_r-t_s|),\qquad
 V=\frac{C_3(M/\log^3M)}{(aM/\log^2M)^2}.
\end{equation}
The arithmetic factors \(\sig(t_r)\sig(t_s)\) in the unnormalized
sieve bound are divided out before applying the source lower bounds.
Using only a pointwise bound on the separation factor would be too
expensive.

\begin{lemma}[Averaged overlap]\label{lem:overlap}
There is \(C_{\mathrm{ov}}>0\) such that, for every sufficiently
large \(M\) and every \(B\subseteq W(M)\) of size
\(b\ge M/\log^4M\),
\begin{equation}\label{eq:overlap}
 \sum_{r\in B}\sum_{\substack{s\in B\\s\ne r}}
       \sig(|t_r-t_s|)
 \le C_{\mathrm{ov}}b^2
       \log\!\left(3+\log\frac{N(M)}b\right).
\end{equation}
\end{lemma}

The formal proof obtains a uniform high-moment bound for the
separation factor and applies H\"older with
\(q=\lceil\log(N(M)/b)\rceil\). In the stated range,
\(q\ll\log\log M\), \(b\le N(M)/2\) for large \(M\), and the
moment estimate costs \(O(\log(q+2))\). This gives precisely the
double logarithm in \eqref{eq:overlap}. Both residue branches and
the omitted diagonal are retained. The lemma is
\lean{exists_uniform_primeTarget_overlap}; it is uniform over
the subset \(B\), which is essential for generated states.

\begin{lemma}[Joint failed-target bound]\label{lem:joint}
Let \(B=G\setminus\cT(S)\), \(b=\card B\), and suppose
\(b\ge M/\log^4M\). If
\begin{equation}\label{eq:absorb}
 8EVC_{\mathrm{ov}}\log\!\left(3+\log\frac{N(M)}b\right)\le1,
\end{equation}
then \(b\le8EU\).
\end{lemma}

\begin{proof}
By \eqref{eq:half-deleted}, each failed target satisfies
\(1\le2\sum_{p\in D}w_r(p)\). Sum over \(B\) and apply
Cauchy--Schwarz:
\[
 b^2\le4E\sum_{p\in D}\left(\sum_{r\in B}w_r(p)\right)^2
 \le4E\sum_{p\in P}\left(\sum_{r\in B}w_r(p)\right)^2.
\]
Expanding the square, \eqref{eq:self}, \eqref{eq:cross}, and
\eqref{eq:overlap} bound the last sum by
\[
 bU+b^2VC_{\mathrm{ov}}\log\!\left(3+\log\frac{N(M)}b\right).
\]
Move the second contribution to the left using \eqref{eq:absorb}.
This gives \(b^2/2\le4EbU\). Since \(b>0\), divide by \(b/2\).
This is a joint estimate; no distribution model for \(D\) was used.
\end{proof}

\begin{proposition}[Row-normalized recurrence]\label{prop:recurrence}
For each \(a>0\), there is \(\eta_a>0\) such that, for every sufficiently
large \(M\), every finite set \(S\), and every exceptional set \(Z\)
satisfying \eqref{eq:supply} on \(W(M)\setminus Z\),
\[
 E\le\eta_a N(M)\quad\Longrightarrow\quad
 \card{\PP(\floor{5M/4})\setminus\cT(S)}
 \le\frac{17}{16}E+\card Z+\frac{8\eta_a}{a}\log M.
\]
\end{proposition}

\begin{proof}
Write \(F=M/\log^4M\) and again let \(b\) count the good new targets
that fail. We prove
\begin{equation}\label{eq:b}
 b\le E/16+(8\eta_a/a)\log M.
\end{equation}
There is nothing to prove when \(b<E/16\). Otherwise
\(b\ge E/16\). Put \(d=E/N(M)\) and \(z=N(M)/b\).
If \(E>0\), then \(0<d\le\eta_a\) and \(1\le z\le16/d\).
The elementary limit
\[
 d\log(3+\log(16/d))\longrightarrow0\qquad(d\downarrow0)
\]
allows \(\eta_a\) to be chosen so small that
\[
 8d(C_3C_{\mathrm{ov}})\log(3+\log z)\le a^2
 \quad(0<d\le\eta_a,\ 1\le z\le16/d).
\]
After substituting \(V\), this is exactly \eqref{eq:absorb}.
Thus, if also \(b\ge F\), Lemma~\ref{lem:joint} gives
\[
 b\le8EU\le (8\eta_a/a)\log M.
\]
When \(E=0\), the half-deleted inequality already forces \(b=0\).

It remains to handle \(0<b<F\). Choose any failed good target.
Its supply and \eqref{eq:half-deleted}, together with
\(\sig(t_r)\ge1\), imply
\(aM/\log^2M\le2E\).
Decrease \(\eta_a\), if necessary, so that
\(64\eta_a\le a\) and \(64\eta_a\le a^2\).
Combining the supply with \(E\le\eta_a M/\log M\) gives
\(a\le2\eta_a\log M\), hence \(a\log^2M\ge32\).
Consequently
\[
 F=\frac{M}{\log^4M}\le
 \frac{aM}{32\log^2M}\le E/16.
\]
This proves \eqref{eq:b} in the small-\(b\) case as well.

An old hole contributes at most \(E\), an exceptional new target at
most \(\card Z\), and a failed good target at most \(b\). Thus
\[
 \card{\PP(\floor{5M/4})\setminus\cT(S)}\le E+\card Z+b.
\]
Use \eqref{eq:b}. This retains all old holes; it does not assume they
are healed.
\end{proof}

\begin{corollary}[Complete one-step estimate]\label{cor:onestep}
There exist \(C,D,\eta_{\mathrm{pow}}>0\) and a scale threshold such
that, for every larger \(M\) and every finite \(S\) with
\(E\le\eta_{\mathrm{pow}}N(M)\),
\begin{equation}\label{eq:onestep}
 \card{\PP(\floor{5M/4})\setminus\cT(S)}
 \le\frac{17}{16}E+CM^{79/80}+D\log M.
\end{equation}
\end{corollary}
\begin{proof}
Use Proposition~\ref{prop:supply} in
Proposition~\ref{prop:recurrence}, and take
\(D=8\eta_{\mathrm{pow}}/a\).
The installed library also provides a proof through its earlier
logarithmic deletion recurrence; either version has exactly the
form \eqref{eq:onestep}.
\end{proof}

\section{Iteration: density zero and the \(79/80\) bound}\label{sec:iteration}

We first record why entry leads to vanishing relative defect.
This intermediate step lets us use independently selected small
density constants without identifying them.

\begin{lemma}[Vanishing defect after entry]\label{lem:zero}
There are \(\eta_{\mathrm{ent}}>0\) and \(M_{\mathrm{ent}}\) such that
\eqref{eq:entry} with these constants implies
\[
 \frac{E_j}{m_j/\log m_j}\longrightarrow0.
\]
\end{lemma}

\begin{proof}
Since \(\gamma<1\), the source exception satisfies
\(CM^\gamma=O(M/\log^4M)\). The logarithmic remainder is also
\(O(M/\log^4M)\). Hence for some fixed \(C_0>0\), the complete
one-step inequality in its applicable density range gives
\[
 E_{j+1}\le\lambda E_j+C_0F_j,\quad
 \lambda=17/16,\quad F_j=m_j/\log^4m_j.
\]
For sufficiently large starting \(M\), both
\(N_j=m_j/\log m_j\) and \(F_j\) grow by at least \(9/8\)
at each step. Choose \(\eta_{\mathrm{ent}}\) within the permitted
density range and then take \(M_{\mathrm{ent}}\) large enough
that \(C_0F_j\le\eta_{\mathrm{ent}}N_j/16\) throughout the
increasing cutoff sequence. Induction gives
\(E_j\le\eta_{\mathrm{ent}}N_j\), so every application of the
recurrence is justified.

Iteration, using \(F_{j+1}\ge(9/8)F_j\), gives
\[
 E_j\le\lambda^jE_0+16C_0F_j.
\]
For example, the inductive step is valid because
\(16C_0(9/8-\lambda)=C_0\).
Divide by \(N_j\). The first term is bounded by
\((17/18)^j E_0/N_0\), and the second is
\(16C_0/\log^3m_j\). Both tend to zero. Floors cause no
difficulty because all the scale comparisons hold beyond a
single fixed threshold.
\end{proof}

\begin{lemma}[Power growth]\label{lem:powergrowth}
For \(M\ge8\), the cutoffs \eqref{eq:cutoffs} satisfy
\[
 m_j\ge M,\qquad
 \left\lfloor\frac{5m_j}{4}\right\rfloor\ge\frac98m_j,\qquad
 m_{j+1}^{\gamma}\ge\frac{13}{12}m_j^\gamma.
\]
\end{lemma}
\begin{proof}
The first assertion follows inductively. For an integer \(m\ge8\),
the floor error in \(5m/4\) is at most \(3/4\), and the stated
\(9/8\) lower bound follows. Finally,
\[
 \left(\frac{13}{12}\right)^4<
 \left(\frac98\right)^3,\qquad
 \gamma=\frac{79}{80}\ge\frac34.
\]
Taking fourth roots and then using monotonicity in the exponent gives
\((9/8)^\gamma\ge13/12\). Raise the cutoff comparison to
the positive power \(\gamma\).
\end{proof}

\begin{lemma}[Eventual recurrence with a growing forcing scale]
\label{lem:abstract}
Suppose \(F_j>0\), \(F_{j+1}\ge(13/12)F_j\) for every \(j\), and,
for all \(j\ge J\),
\[
 E_{j+1}\le(17/16)E_j+C_1F_j.
\]
Then there is \(K>0\) such that \(E_j\le KF_j\) for every \(j\ge J\).
One may take
\[
 K=\max\{1,E_J/F_J,48C_1\}.
\]
\end{lemma}
\begin{proof}
At \(j=J\) this is the definition of \(K\). Since
\(13/12-17/16=1/48\), we have \(C_1\le K/48\). Thus
\[
 E_{j+1}\le\left(\frac{17}{16}K+C_1\right)F_j
 \le\frac{13}{12}KF_j\le KF_{j+1}.
\]
Induction proves the result. No recurrence before \(J\) is required.
\end{proof}

\begin{proof}[Proof of Theorem~\ref{thm:main}: the actual generations]
Choose entry constants from Lemma~\ref{lem:zero}, and increase the
threshold to at least \(8\). Fix \(M,S\) satisfying entry.
Apply Corollary~\ref{cor:onestep}, whose density constant
\(\eta_{\mathrm{pow}}\) need not equal the entry constant.
Lemma~\ref{lem:zero} implies \(E_j\le\eta_{\mathrm{pow}}N(m_j)\)
eventually. Also \(m_j\to\infty\). Thus the complete power recurrence
\eqref{eq:onestep} applies to the actual set \(S_j\) for every
sufficiently large \(j\).

For \(X>0\), the elementary inequality \(\log X\le2X^\gamma\)
follows from \(\log(X^\gamma)\le X^\gamma-1\).
Hence eventually
\[
 E_{j+1}\le\frac{17}{16}E_j+(C+2D)m_j^\gamma.
\]
Set \(F_j=m_j^\gamma\). Lemmas~\ref{lem:powergrowth}
and \ref{lem:abstract} prove \eqref{eq:stage}.
This step keeps the seed fixed while the generation grows; a different
seed is not selected at each scale.
\end{proof}

\begin{proof}[Proof of Theorem~\ref{thm:main}: arbitrary cutoffs]
Fix \(K,J\) from \eqref{eq:stage}. Given \(X>m_J\), choose the first
\(j\) such that \(X\le m_j\). Monotonicity gives \(j\ge J\), while
minimality gives \(m_{j-1}<X\) and
\[
 m_j=\floor{5m_{j-1}/4}\le5X/4.
\]
Every permanently missing prime at most \(X\) belongs to
\(\PP(m_j)\setminus S_j\). Therefore
\[
 \card{\PP(X)\setminus\Cl(S)}
 \le E_j\le Km_j^\gamma\le K(5/4)^\gamma X^\gamma.
\]
Take \(K'=K(5/4)^\gamma\). This proves \eqref{eq:closure}
for all sufficiently large integer \(X\).
\end{proof}

\section{Birth layers, exact cancellation, and what remains open}
\label{sec:birth}

The previous results begin with a nearly complete prime prefix.
For small-seed growth one must show that a significant part of the
representation mass produces \emph{distinct new owners}.
The development contains exact tools for that question, but not the
needed positive lower bound.

Let \(A,S\) be finite prime sets with \(\cT(A)\subseteq S\), and let
\(B\) be any finite set of prime targets \(r\ge3\) outside \(S\).
Then
\begin{equation}\label{eq:cancel}
 R_A(r)=0\qquad(r\in B).
\end{equation}
Otherwise a representation from \(A\) would put \(r\) in
\(\cT(A)\subseteq S\). Thus for every real \(\beta\),
\[
 \frac{R_S(r)-\beta R_A(r)}{\sig(t_r)}
   =\frac{R_S(r)}{\sig(t_r)}\qquad(r\in B).
\]
This is exact birth-layer cancellation, not an independence heuristic.

If \(S\subseteq\PP(M)\) and \(M\ge5^8\), the two-prime upper sieve gives
\[
 0\le w_r:=R_S(r)/\sig(t_r)\le
 H(M):=C_2M/\log^2M.
\]
Writing \(W_B=\sum_{r\in B}w_r\), one obtains
\begin{equation}\label{eq:energy}
 \sum_{r\in B}
 \left(\frac{R_S(r)-\beta R_A(r)}{\sig(t_r)}\right)^2
 \le H(M)W_B.
\end{equation}
Indeed \eqref{eq:cancel} reduces each summand to \(w_r^2\le H(M)w_r\).
Since a positive \(w_r\) produces \(r\), this also yields
\begin{equation}\label{eq:newcount}
 W_B\le H(M)\card{B\cap(\cT(S)\setminus S)}.
\end{equation}
For actual consecutive generations take \(A=S_{j-1}\) and \(S=S_j\).

These inequalities control the energy in terms of the fresh mass.
They do not prove that \(W_B\) is positive or grows along the orbit
from \(\{2\}\). In particular, an upper bound such as
\eqref{eq:energy} cannot be substituted for the lower bound needed in
\eqref{eq:newcount}. The open small-seed problem is to obtain a
usable joint newness/collision estimate that supplies enough fresh
mass to reach a certified density region.

\subsection{The role of the prime \(2\)}

For prime inputs, \(2(p+q)-1\) is odd and at least \(7\), so its largest
prime factor is odd. Thus \(2\) is never created and must be included
if the intended closure is to contain every prime. Monotonicity of
\(\cT\) gives the seed-inclusion reductions used in the implementation.
The diagonal from \(p\) produces \(P^{+}(4p-1)\); it is available at
every generation. For instance, \(\{3\}\) is not closed: the first
diagonal produces \(11\), then the diagonal from \(11\) produces \(43\).

\section{Finite verification and quantitative scope}\label{sec:finite}

The repository contains a kernel-checked finite certificate
\[
 \PP(10000)\subseteq S_9\qquad\text{for }S_0=\{2\}.
\]
Its public owner is
\lean{inputPrimePrefix_10000_subset_fullOrbit_two_9}.
This proves the displayed finite inclusion. It does not bridge the
entry threshold in Theorem~\ref{thm:main}.

There is also a separate explicit density-preservation theorem.
Writing \(\eta_{\mathrm{cert}}\) for the installed certified fraction,
one has \(\eta_{\mathrm{cert}}\ge2^{-600}\), and
\[
 M\ge2^{49152},\quad
 \card{\PP(M)\setminus S}\le\eta_{\mathrm{cert}}\frac M{\log M}
 \ \Longrightarrow\
 \card{\PP(\floor{7M/6})\setminus\cT(S)}
 \le\eta_{\mathrm{cert}}\frac{\floor{7M/6}}{\log\floor{7M/6}}.
\]
This statement uses expansion factor \(7/6\), not \(5/4\).
It provides a conservative explicit preservation range, not a numerical
value for the eventual threshold of the new \(79/80\) theorem.
The source cutoff \(2^{44}\), the preservation cutoff \(2^{49152}\),
and the finite certificate at \(10000\) have different consumers.

The strongest completed dynamical conclusion is therefore the
power-saving theorem for sufficiently dense finite seeds.
A precise next target is to prove that some actual generation from
\(\{2\}\) satisfies a fully certified entry condition, or to derive a
growth-and-density bootstrap that establishes such entry. Neither a
finite diagnostic nor the density result for all prime-pair inputs
settles that target. Dense entry is an intermediate milestone toward
Conjectures~\ref{conj:two} and \ref{conj:odd}; it would not by itself
remove every exceptional prime.

\subsection{Affine and higher-arity generation questions}
\label{sec:extensions}

The broader proposed programme concerns fixed integers
\[
 a\ge1,\qquad b\in\mathbb Z,\qquad m\ge2.
\]
For a nonempty finite prime seed \(S\), define
\[
 \begin{aligned}
 \cT_{a,b,m}(S)=S\cup\bigl\{
 &P^{+}(\lvert a(p_1+\cdots+p_m)+b\rvert):
       (p_1,\ldots,p_m)\in S^m,\\
 &\lvert a(p_1+\cdots+p_m)+b\rvert>1
 \bigr\}.
 \end{aligned}
\]
Let \(\Cl_{a,b,m}(S)\) be its increasing iterated closure. Tuples may
repeat inputs; permutations and repeated representations do not create
multiple copies of an output prime. The absolute value and the exclusion
of \(0,\pm1\) specify the convention when \(b\) is negative.
The present operator is the case \((a,b,m)=(2,-1,2)\).

We conjecture that complete prime-generation phenomena extend to suitable
affine parameters and to sums of more than two prime inputs, with
appropriate restrictions on the seed and target primes. A precise
classification of those parameters and restrictions remains open.
This is a proposed extension, not a theorem covering all triples.

Some necessary restrictions are immediate. If a prime \(\ell\mid a\)
does not divide \(b\), then every pre-factorization value is nonzero
modulo \(\ell\). Consequently \(\ell\) can belong to the closure only
if it is already in the seed. For example, \(a=2,b=-1\) excludes the
creation of \(2\) at every arity. When all available inputs are odd,
the value before factorization has parity \(am+b\); this must also
be retained when formulating a candidate generation statement.
A fixed divisor, on the other hand, is not necessarily the largest
prime factor, so divisibility alone does not identify the output support.

Congruence restrictions do not exhaust the possible degeneracies.
For \((a,b,m)=(1,0,2)\), every prime singleton is closed, since
\(P^{+}(p+p)=p\). Thus one cannot simply assert universal generation
for every affine map after mentioning modulus caveats.
The general problem is to determine which non-degenerate maps and seeds
generate every target prime compatible with their arithmetic restrictions,
and then prove individual reachability for those cases.

The canonical-owner reduction has a direct higher-arity form. For a
target prime \(r\), it is sufficient to find \(h\ge1\), all of whose prime
factors are at most \(r\), and available inputs satisfying
\[
 a(p_1+\cdots+p_m)+b=hr.
\]
Then the largest prime factor is exactly \(r\). If all inputs are odd,
the sum must be congruent to \(m\) modulo \(2\). In the subfamily
\((a,b)=(2,-1)\), this requires
\[
 hr\equiv2m-1\pmod4.
\]
For every odd prime \(r\), one of \(h=1\) and \(h=3\) satisfies this
congruence, and its prime factors are at most \(r\). This removes the
elementary parity obstruction to that owner representation; it does
not establish the required additive representation from a generated
state. The higher-arity direction therefore preserves the distinction
between an arithmetic compatibility check and an actual reachability proof.

\appendix
\section{Formal statements and verification map}\label{app:formal}

The maintained proof uses Windows Lean \(4.34.0\), the repository's
pinned Mathlib environment, and released dependency artifacts.
Dependencies are reused from their existing caches. The main theorem
has no RH, Goldbach, mixing, or small-seed-entry hypothesis.
Its substantive seed hypothesis is precisely \eqref{eq:entry}.
The analytic facts used to discharge supply and overlap belong to the
transitive proof dependency closure.

The following names are relative to
\lean{PrimeFactorOscillations.LPF}, except for the final public consumer.
Long implementation paths are listed by module stem to keep the table readable.

\small
\begin{longtable}{@{}>{\raggedright\arraybackslash}p{0.30\textwidth}>{\raggedright\arraybackslash}p{0.65\textwidth}@{}}
\toprule
Mathematical item & Maintained owner or declaration\\
\midrule
\endhead
Full dynamics with diagonals &
\lean{LPFDynamics}: \lean{fullStep}, \lean{fullOrbit}. The older
\lean{step} and \lean{orbit} definitions exclude diagonals and are
not the dynamics of this paper.\\[3pt]
Canonical ownership &
\lean{LPFJointTransfer}: \lean{primeTarget_lift},
\lean{primeTarget_owner}, \lean{primeTarget_injective_window}.\\[3pt]
Uniform character prefix &
\lean{LPFCharacterPrefix}: namespace \lean{QRH},
\lean{exists_all_characterChebyshevSum_power_bound}.\\[3pt]
Four-packet source identity and budget &
\lean{LPFFixedSourceBudget}: \lean{lpf_fixed_source_four_packet_budget};
\lean{LPFSourcePower}: \lean{lpf_fixed_source_power_bound}.\\[3pt]
Positive representation supply &
\lean{LPFWideSource}:
\lean{exists_wide_primeTarget_supply_power_explicit}.\\[3pt]
Subset-uniform overlap &
\lean{LPFPrimeTargetOverlap}:
\lean{exists_uniform_primeTarget_overlap}.\\[3pt]
Row normalization and deletion &
\lean{LPFRowNormalization}:
\lean{primeTargetWeight_self_exact},
\lean{exists_dense_fullStep_prefix_row_normalized_recurrence}.\\[3pt]
Complete power/log recurrence &
\lean{LPFPositiveSupply}:
\lean{exists_fullStep_dense_power_log_bound_from_wide_source}.\\[3pt]
Vanishing relative defect &
\lean{LPFDenseDensity}:
\lean{exists_fullOrbit_dense_density_zero}, with its representation
premise discharged by \lean{exists_wide_primeTarget_supply_log}.\\[3pt]
Actual-generation \(79/80\) theorem &
\lean{LPFDensePower}:
\lean{exists_fullOrbit_dense_power_bound}.\\[3pt]
All-cutoff closure theorem &
\lean{LPFDensePower}:
\lean{fullOrbitMissingPrimes},
\lean{exists_fullOrbit_dense_closure_power_bound}.\\[3pt]
Public implementation consumer &
\lean{Solution.lean}:
\lean{PrimeFactorOscillations.publicLPFDenseClosurePowerBound}.\\[3pt]
Fresh-owner energy and mass &
\lean{LPFWeightedCanonicalEnergy}:
\lean{fresh_birth_energy_upper}, \lean{fresh_mass_bound}.\\[3pt]
Finite seed certificate &
\lean{LPFFiniteEntry.Result}:
\lean{inputPrimePrefix_10000_subset_fullOrbit_two_9}.\\[3pt]
Separate explicit preservation &
\lean{LPFNumericalPreservation}:
\lean{lpf_density_preservation_numerical_cutoff}.\\
\bottomrule
\end{longtable}
\normalsize

The closure statement is formalized over integer cutoffs. In readable
quantifier form it says
\[
\begin{aligned}
 &\exists\eta>0\ \ \forall M\text{ sufficiently large}\ \ \forall S,\\
 &\quad \card{\PP(M)\setminus S}\le\eta M/\log M\\
 &\qquad\Longrightarrow\
   \exists K>0\ \ \forall X\text{ sufficiently large},\\
 &\hspace{42mm}\card{\PP(X)\setminus\Cl(S)}\le KX^{79/80}.
\end{aligned}
\]
The implementation represents the final missing set by filtering the finite
prime prefix with the condition that membership fails at every generation.
No infinite-set cardinality is used.

The formal proof of the new power-saving conclusion proceeds in the same
order as Section~\ref{sec:iteration}: cutoff power growth, an eventual
recurrence lemma, the transfer from vanishing density to the power
recurrence, and first-crossing interpolation. These are proofs, not axioms
encoding the conclusion. The declaration audit checks their complete
axiom sets; standard foundational axioms such as propositional
extensionality, classical choice, and quotient soundness are distinguished
from any additional analytic assumption. The build and audit commands,
artifact pins, and saved results are supplied in the companion
\href{https://github.com/kimihiro64/prime-factor-oscillations}{repository}.

\section{What the result does and does not transfer}

The result concerns an actual iterated closure once finite entry is
given. It is stronger than a limiting divisibility law for the full
prime-pair ensemble: such a law alone does not describe a dynamically
generated input set. Conversely, it is weaker than all-prime
reachability. A power-saving exceptional set may still contain infinitely
many primes.

No criterion for the Riemann hypothesis or for bounded prime gaps is
derived from this dynamical theorem. The companion paper on
prime-factor densities treats those separate questions with different
quantifiers and different signed inequalities. Sharing analytic tools
does not supply an implication between the two sets of conclusions.

\clearpage
\begin{thebibliography}{9}
\bibitem{BC}
G.~Back and M.~Caragiu,
\emph{The greatest prime factor and recurrent sequences},
Fibonacci Quarterly \textbf{48} (2010), no.~4, 358--362.
\url{https://fq.math.ca/Papers1/48-4/Back_Caragiu.pdf}.

\bibitem{CV}
M.~Caragiu and P.~A.~Vicol,
\emph{Prime magmas and a cyclicity conjecture},
JP Journal of Algebra, Number Theory and Applications
\textbf{38} (2016), no.~2, 129--143.
\href{https://doi.org/10.17654/NT038020129}{doi:10.17654/NT038020129}.

\bibitem{RR}
O.~Ramar\'e and I.~Z.~Ruzsa,
\emph{Additive properties of dense subsets of sifted sequences},
Journal de Th\'eorie des Nombres de Bordeaux
\textbf{13} (2001), no.~2, 559--581.
\href{https://doi.org/10.5802/jtnb.338}{doi:10.5802/jtnb.338}.

\bibitem{GT}
B.~Green and T.~Tao,
\emph{Restriction theory of the Selberg sieve, with applications},
Journal de Th\'eorie des Nombres de Bordeaux
\textbf{18} (2006), no.~1, 147--182.
\url{https://jtnb.centre-mersenne.org/article/JTNB_2006__18_1_147_0.pdf}.

\bibitem{AS}
A.~Alsetri and X.~Shao,
\emph{Density versions of the binary Goldbach problem},
Acta Arithmetica \textbf{218} (2025), no.~3, 285--295.
Preprint arXiv:2405.18576v2, 19 September 2024.
\url{https://arxiv.org/abs/2405.18576v2}.

\bibitem{OAI}
OpenAI,
\emph{The Quasi-Riemann Hypothesis: A Zero-Free Half-Plane
\(\operatorname{Re}(s)>7/8\)},
preprint, 30 September 2026, with companion Lean development.
\url{https://github.com/openai/math/tree/main/preprints/The-Quasi-Riemann-Hypothesis-September-30-2026}.
The referenced nonvanishing source is also available at revision
\lean{adc7f1241b42e322a6451854ab7e4b4c146bf78a} of that repository.

\bibitem{Repo}
J.~Whidden,
\emph{Prime Factor Oscillations}, companion source and formal proofs,
2026.
\url{https://github.com/kimihiro64/prime-factor-oscillations}.
\end{thebibliography}
\end{document}
